\documentclass[10pt,a4paper]{amsart}
\usepackage[margin=19mm]{geometry}
\usepackage{amsmath,amssymb,mathtools}
\usepackage[scr=rsfs,cal=euler]{mathalfa}
\usepackage{xfrac}
\usepackage[unicode,hidelinks,bookmarks=false]{hyperref}
\newcommand{\ie}{i.e.\ }
\newcommand{\cf}{cf.\ }
\newcommand{\resp}{respectively\ }
\newcommand{\Subsection}{\subsection}
\providecommand{\nobreakdash}{\nobreak\hbox{-}}
\setlength{\emergencystretch}{2em}
\multlinegap0pt
\usepackage{bm}
\usepackage{bbm}
\usepackage{dsfont}
\usepackage{xspace}
\usepackage{cancel}
\usepackage{mathtools}
\usepackage{amsthm}
\makeatletter
\@ifundefined{theorem}{\newtheorem{theorem}{Theorem}}{}
\makeatother
\newcommand{\ra}{\rightarrow}
\title[Extended Kato inequalities for conformal operators]{Extended Kato inequalities for conformal operators}
\author{Daniel Cibotaru, Matheus Vieira}
\date{Source: arXiv:2410.15520v2 (CC BY 4.0)}
\begin{document}
\maketitle

\noindent\emph{Source and licence.} Extended Kato inequalities for conformal operators. Version arXiv:2410.15520v2 (CC BY 4.0), distributed under the Creative Commons Attribution 4.0 International licence, as recorded in the arXiv licence metadata for this exact version and shown on the abstract page https://arxiv.org/abs/2410.15520v2. Full rights notice: licenses/arxiv-2410.15520-CC-BY-4.0.txt.\par\smallskip

\noindent\textit{Editorial scope.} Counted unit: the Theorem whose source label is \texttt{Dinjelliptic} in the article (source0.tex lines 584--589, source label \texttt{Dinjelliptic}), delivered together with the article's own complete original proof of it (source0.tex lines 591--612, one \texttt{proof} environment of 22 lines). No mathematics is written, repaired or re-derived: statement and proof are the article's own text line for line. Required context retained uncounted: none. Every definition, display and label of those lines is retained unchanged. Omitted from this package and not redistributed: 1--583, 590--590, 613--1090. Nothing inside a retained excerpt is omitted or altered: every retained body is delivered exactly as the article's lines are written, so no omission receipt of any kind accompanies this package. Citations: the retained excerpts carry no citation command at all, so no bibliography entry of the article is transcribed and no reference is redistributed. Reference adaptation: none was needed and none was made; every reference inside the delivered unit resolves inside it, so no unresolved reference or citation is produced. Printed slips kept literal: no printed word, symbol, hypothesis or number of the counted statement or of its proof was altered. Layout and class: the article's own class is \texttt{amsart} with packages including amssymb, amsmath, latexsym, amscd, amsfonts, pb-diagram, xy, eucal, amsthm, epsfig, url, pdfsync, bbm, hyperref; the wrapper is the lane's pinned \texttt{amsart} editorial wrapper at 10pt on A4, which restates the theorem-like environments the retained text needs and adds the article's own macro definitions that the retained text needs (\texttt{\textbackslash ra}), with extra wrapper packages (bm, bbm, dsfont, xspace, cancel, mathtools). The delivered numbering is the unit's own. Assets: no retained span contains a figure environment, a graphics-inclusion command, a PDF-page inclusion, a table or a TikZ picture; no image file is copied into this package, the package declares no asset dependency and no image file is redistributed with it. Licence: version arXiv:2410.15520v2 of this article is distributed under the Creative Commons Attribution 4.0 International licence, as recorded in the arXiv licence metadata for this exact version and shown on the abstract page https://arxiv.org/abs/2410.15520v2. A full rights notice is delivered at licenses/arxiv-2410.15520-CC-BY-4.0.txt. Characters: every retained excerpt is pure ASCII, so the delivered engine, XeLaTeX with its default Latin Modern text font, reports no missing character.
\begin{theorem}\label{Dinjelliptic} Let $D:\Gamma(E)\ra\Gamma(F)$ be a conformal injectively elliptic operator  between Riemannian vector bundles $E,F \ra M$ with ellipticity constant $\epsilon$ and conformity factor $\rho$. Let $\phi$ be a section of $E$ and let $x$ be a point in $M$ such that $\phi(x)\neq0$. Then at $x$ the following pointwise inequality holds
\begin{equation}\label{fexK} \sqrt{\rho^ 2- \epsilon}|\nabla \phi|+ |D \phi|\geq  \rho|d|\phi||.
\end{equation}
In particular, if $(D\phi)_x=0$ then the following refined Kato inequality holds
\[|\nabla \phi|^ 2\geq \left(1+\frac{\epsilon}{\rho^2-\epsilon}\right)|d|\phi||^ 2. \]
\end{theorem}

\begin{proof}
If $|d|\phi||(x) = 0$ the result is trivial, so by continuity we can assume that $|d|\phi|| \ne 0$ in a neighborhood of $x$. We take the one-form $\xi_0:=\frac{d|\phi|^2}{|d|\phi|^2|}$, which is defined in a neighborhood of $x$ and has norm $1$. Then
\begin{equation}\label{equ5} |d|\phi||\cdot|\phi|=\langle \nabla_{\xi_0^{\sharp}}\phi,\phi\rangle=\langle \nabla \phi,\xi_0\otimes \phi \rangle ,
\end{equation}
where $(\cdot)^{\sharp}$ is the musical isomorphism. The justification is immediate:
\begin{equation}\label{mine} 2|\phi|\cdot|d|\phi||=|d|\phi|^2|=\langle d|\phi|^2,\xi_0\rangle=\xi_0^{\sharp}|\phi|^2=2\langle \nabla_{\xi_0^{\sharp}}\phi,\phi \rangle=2\langle \nabla \phi,\xi_0\otimes \phi \rangle.
\end{equation}
The operator $\frac{1}{\rho}P:T^ *M\otimes E\ra F$ is an orthogonal projection. Then there exists a projection $\pi^{F^ {\perp}}:T^ *M\otimes E\ra F^ {\perp}$ onto the orthogonal complement of $P^ *(F)$, and so
\[T^ *M\otimes E=P^ *(F)\oplus F^ {\perp}.\]
Correspondingly, we have \[\nabla \phi=\frac{1}{\rho}D\phi + \pi^{F^{\perp}} (\nabla \phi),\]
and so 
\begin{equation}\label{ineq1}\langle \nabla \phi,\xi_0\otimes \phi \rangle=\frac{1}{\rho}\langle D \phi ,\xi_0\otimes \phi\rangle+\langle \nabla \phi,\pi^{F^ {\perp}}(\xi_0\otimes \phi)\rangle.\end{equation}
We have:
\[\frac{1}{\rho^ 2}|P(\xi_0\otimes \phi)|^ 2+|\pi^ {F^ {\perp}}(\xi_0\otimes \phi)|^ 2=|\xi_0\otimes \phi|^ 2=|\phi|^ 2.\]
But the injective ellipticity condition implies that
\[|P(\xi_0\otimes \phi)|^ 2\geq \epsilon |\xi_0|^ 2|\phi|^ 2=\epsilon|\phi|^2, \]
and therefore
\[ \left(1-\frac{\epsilon}{\rho^ 2}\right)|\phi|^2 \geq |\pi^ {F^ {\perp}}(\xi_0\otimes \phi)|^2.\]
Using (\ref{ineq1}) and the Cauchy-Schwarz inequality we get
\begin{equation}\label{equa6}  \sqrt{1-\frac{\epsilon}{\rho^ 2}}|\nabla \phi|\cdot|\phi|+\frac{1}{\rho}|D(\phi)|\cdot|\phi|\geq |\langle \nabla \phi,\xi_0\otimes \phi \rangle|.\end{equation}
Now (\ref{equa6}) with (\ref{equ5}) gives (\ref{fexK}).
\end{proof}

\end{document}
